\begin{abstract}
Temporal graph-fraud results are usually reported for a fixed split, graph, and metric, although deployment also fixes when labels arrive, which edges are visible, how the graph is constructed, how many alerts can be reviewed, and what computation is feasible. A model ranking without those conditions has no stable operational meaning. We introduce FraudShiftBench, which represents these conditions as a deployment contract and binds each scoped conclusion to a typed evidence unit. Its support relation checks cell completeness, prediction provenance, paired uncertainty, and resource status before a claim can enter the benchmark. We instantiate the framework with ten-seed evidence on two real financial graphs and four synthetic IBM AML-Data regimes. On Elliptic, changing only strict versus isolated graph access changes the AUPRC leader from GraphSAGE to GCN; their paired mean changes are $-0.132$ and $+0.094$, respectively, while the feature-only MLP is invariant. In the IBM baseline grid, histogram gradient boosting has the highest mean AUPRC in all eight variant--protocol cells. Across the broader feasible IBM study, however, AUPRC and fixed-threshold F1 select different configurations in 11 of 16 contexts. Within the graph-configuration grid, edge-conditioned GINE leads AUPRC in all four measured Small contexts, but its Medium cells remain outside the ranking after T4 memory exhaustion. Fourteen claim-mutation and evidence-ablation tests produce the specified status transitions, including rejection of incomplete, prediction-missing, and construct-invalid claims. The resulting benchmark is a conditional map of protocol, architecture, construction, metric, scale, and feasibility rather than a global leaderboard. The released artifact contains prediction manifests, evidence locks, executable validators, deterministic analyses, and an explicit record of exclusions and resource boundaries.
\end{abstract}

\begin{IEEEkeywords}
Graph fraud detection, anti-money laundering, temporal shift, evaluation protocols, benchmark validity, reproducibility.
\end{IEEEkeywords}
